\NormalFont\ShortTitle{2 Petru}

{\MT A doua scrisoare a lui Petru


\par }\ChapOne{1}{\PP \VerseOne{1}Simon Petru, slujitor și apostol al lui Isus Hristos, către cei ce au dobândit o credință la fel de prețioasă ca și noi în dreptatea Dumnezeului și Mântuitorului nostru Isus Hristos:

\VS{2}Harul să vă fie vouă și pacea să se înmulțească în cunoașterea lui Dumnezeu și a lui Isus, Domnul nostru,

\VS{3}având în vedere că puterea sa divină ne-a acordat toate lucrurile care țin de viață și de evlavie, prin cunoașterea celui care ne-a chemat prin propria sa glorie și virtute,

\VS{4}prin care ne-a acordat prețioasele și nespus de mari făgăduințe, pentru ca, prin acestea, să deveniți părtași ai naturii divine, după ce ați scăpat de corupția care este în lume prin pofte.

\VS{5}Da, și tocmai de aceea, adăugând din partea voastră toată sârguința, în credința voastră supliniți excelența morală; și în excelența morală, cunoașterea;

\VS{6}și în cunoaștere, stăpânirea de sine; și în stăpânirea de sine, perseverența; și în perseverență, evlavia;

\VS{7}și în evlavie, afecțiunea frățească; și în afecțiunea frățească, dragostea.

\VS{8}Căci, dacă aceste lucruri sunt ale voastre și abundă, ele vă fac să nu fiți leneși și să nu rămâneți fără rod în cunoașterea Domnului nostru Isus Hristos.

\VS{9}Căci cel căruia îi lipsesc aceste lucruri este orb, văzând doar ceea ce este aproape, uitând curățirea de vechile sale păcate.

\VS{10}De aceea, fraților, fiți mai sârguincioși să vă asigurați chemarea și alegerea voastră. Căci, dacă faceți aceste lucruri, nu vă veți poticni niciodată.

\VS{11}Căci astfel veți avea din belșug intrarea în Împărăția veșnică a Domnului și Mântuitorului nostru, Isus Hristos.


\par }{\PP \VS{12}De aceea nu voi neglija să vă aduc aminte de aceste lucruri, deși le cunoașteți și sunteți întăriți în adevărul prezent.

\VS{13}Mi se pare drept, atâta timp cât sunt în cortul acesta, să vă stârnesc, aducându-vă aminte,

\VS{14}știind că desființarea cortului meu vine repede, așa cum mi-a arătat Domnul nostru Isus Hristos.

\VS{15}Da, mă voi strădui ca voi să vă puteți aminti mereu aceste lucruri, chiar și după plecarea mea.


\par }{\PP \VS{16}Fiindcă nu am urmărit minciuni bine gândite, când v-am făcut cunoscută puterea și venirea Domnului nostru Isus Hristos, ci am fost martori oculari ai măreției Lui.

\VS{17}Căci el a primit de la Dumnezeu Tatăl onoare și glorie, atunci când i-a venit glasul din slava maiestuoasă: “Acesta este Fiul Meu preaiubit, în care Îmi găsesc plăcerea.”

\VS{18}Noi am auzit acest glas venind din cer, când eram cu el pe muntele sfânt.


\par }{\PP \VS{19}Noi avem cuvântul profeției, care este mai sigur; și bine faceți că-l luați în seamă ca pe o lampă care strălucește într-un loc întunecos, până ce se va face ziuă și va răsări steaua dimineții în inimile voastre,

\VS{20}știind mai întâi că nici o profeție a Scripturii nu este de neînțeles.

\VS{21}Căci nicio profeție nu a venit vreodată prin voința omului, ci au vorbit bărbați sfinți ai lui Dumnezeu, mișcați de Duhul Sfânt.



\par }\Chap{2}{\PP \VerseOne{1}Dar în popor au apărut și prooroci mincinoși, cum vor fi și printre voi învățători mincinoși, care vor aduce în ascuns erezii nimicitoare, care vor nega chiar pe Stăpânul care i-a cumpărat și care vor atrage asupra lor o grabnică pieire.

\VS{2}Mulți vor urma căile lor imorale și, ca urmare, calea adevărului va fi defăimată.

\VS{3}În lăcomie vă vor exploata cu vorbe înșelătoare, a căror sentință, acum din vechime, nu zăbovește, și distrugerea lor nu va adormi.


\par }{\PP \VS{4}Căci, dacă Dumnezeu nu a cruțat pe îngeri când au păcătuit, ci i-a aruncat în Tartar și i-a trimis în gropile întunericului, ca să fie rezervați pentru judecată;

\VS{5}și nu a cruțat lumea antică, ci l-a păstrat pe Noe împreună cu alți șapte, un propovăduitor al dreptății, atunci când a adus un potop peste lumea celor nelegiuiți,

\VS{6}și, transformând în cenușă cetățile Sodoma și Gomora, le-a condamnat la pieire, după ce le-a dat ca exemplu celor care vor trăi în mod nelegiuit,

\VS{7}și l-a eliberat pe neprihănitul Lot, care era foarte chinuit de viața plină de pofte a celor nelegiuiți

\VS{8}(căci acel om drept care locuia în mijlocul lor era chinuit în sufletul său drept din zi în zi văzând și auzind fapte nelegiuite),

\VS{9}atunci Domnul știe cum să-i scape pe cei evlavioși din ispită și cum să-i țină pe cei nelegiuiți sub pedeapsă pentru ziua judecății,

\VS{10}dar mai ales pe cei care umblă după trup în pofta desfrânării și disprețuiesc autoritatea. Îndrăzneți, plini de voință proprie, ei nu se tem să vorbească de rău demnitarii,

\VS{11}în timp ce îngerii, deși sunt mai mari în putere și în tărie, nu aduc o judecată calomnioasă împotriva lor înaintea Domnului.

\VS{12}Dar aceștia, ca niște creaturi fără rațiune, născuți ca animale naturale pentru a fi luate și distruse, vorbind de rău în chestiuni despre care nu știu nimic, vor fi cu siguranță distruși în distrugerea lor,

\VS{13}primind plata nedreptății; oameni care socotesc că este o plăcere să se distreze în timpul zilei, cu pete și defecte, bucurându-se de înșelăciunea lor în timp ce petrec cu voi;

\VS{14}având ochii plini de adulter și care nu pot înceta să păcătuiască, ademenind suflete neliniștite, având o inimă antrenată în lăcomie, copii blestemați!

\VS{15}Abandonând calea cea dreaptă, s-au rătăcit, urmând calea lui Balaam, fiul lui Beor, care iubea plata nedreptății;

\VS{16}dar el a fost mustrat pentru neascultarea sa. Un măgar fără glas a vorbit cu voce de om și a oprit nebunia profetului.


\par }{\PP \VS{17}Aceștia sunt fântâni fără apă, nori împinși de furtună, cărora li s-a rezervat pentru totdeauna întunericul întunericului.

\VS{18}Căci, rostind vorbe mari și umflate de goliciune, ei ademenesc în poftele cărnii, prin desfrânare, pe cei care, într-adevăr, scapă de cei care trăiesc în rătăcire;

\VS{19}promițându-le libertate, în timp ce ei înșiși sunt robi ai corupției; căci omul este adus în robie de cel care îl biruiește.


\par }{\PP \VS{20}Căci dacă, după ce au scăpat de spurcăciunea lumii prin cunoașterea Domnului și Mântuitorului Isus Hristos, se vor încurca iarăși în ea și vor fi biruiți, atunci starea de pe urmă va fi mai rea pentru ei decât cea dintâi.

\VS{21}Căci ar fi mai bine pentru ei să nu fi cunoscut calea dreptății, decât ca, după ce au cunoscut-o, să se întoarcă de la porunca sfântă care le-a fost dată.

\VS{22}Dar li s-a întâmplat ceea ce se spune în proverbul adevărat: “Câinele se întoarce la propria vomă” și “scroafa care s-a spălat se tăvălește în noroi”.



\par }\Chap{3}{\PP \VerseOne{1}Aceasta este acum, iubiților, a doua scrisoare pe care v-am adresat-o. Și în amândouă vă trezesc mintea voastră sinceră, amintindu-vă

\VS{2}să vă aduceți aminte de cuvintele care au fost rostite mai înainte de sfinții prooroci și de porunca noastră, a apostolilor Domnului și Mântuitorului,

\VS{3}știind mai întâi că în zilele din urmă vor veni batjocoritori, umblând după poftelelor

\VS{4}și zicând: “Unde este făgăduința venirii Lui?”. Căci, din ziua în care părinții au adormit, toate lucrurile continuă așa cum au fost de la începutul creației”.

\VS{5}Căci ei uită cu bună știință că au existat ceruri din vechime și un pământ format din apă și în mijlocul apei prin cuvântul lui Dumnezeu,

\VS{6}prin care lumea care exista atunci, fiind inundată de apă, a pierit.

\VS{7}Dar cerurile care există acum și pământul, prin același cuvânt, au fost depozitate pentru foc, fiind rezervate pentru ziua judecății și a distrugerii oamenilor nelegiuiți.


\par }{\PP \VS{8}Dar nu uitați, iubiților, că o zi este pentru Domnul ca o mie de ani, și o mie de ani ca o zi.

\VS{9}Domnul nu este încet în ceea ce privește făgăduința Sa, cum socotesc unii că este încetineala, ci este răbdător cu noi, nevrând ca nimeni să piară, ci ca toți să vină la pocăință.

\VS{10}Dar ziua Domnului va veni ca un hoț în noapte, în care cerurile vor trece cu zgomot mare, iar elementele se vor dizolva cu ardoare, iar pământul și lucrările care sunt în el vor fi arse.

\VS{11}Prin urmare, de vreme ce toate aceste lucruri vor fi distruse astfel, ce fel de oameni ar trebui să fiți voi în viață sfântă și evlavie,

\VS{12}așteptând și dorind cu ardoare venirea zilei lui Dumnezeu, care va face să se dizolve cerurile arzând, iar elementele se vor topi cu căldură arzătoare?

\VS{13}Ci, potrivit promisiunii Lui, așteptăm ceruri noi și un pământ nou, în care sălășluiește dreptatea.


\par }{\PP \VS{14}De aceea, preaiubiților, întrucât așteptați aceste lucruri, sârguiți-vă să fiți găsiți în pace, fără cusur și fără vină înaintea Lui.

\VS{15}Priviți răbdarea Domnului nostru ca pe o mântuire; așa cum v-a scris și iubitul nostru frate Pavel, după înțelepciunea care i-a fost dată,

\VS{16}ca și în toate scrisorile sale, vorbind în ele despre aceste lucruri. În acestea, sunt unele lucruri greu de înțeles, pe care cei neștiutori și nestatornici le răstălmăcesc, așa cum fac și cu celelalte Scripturi, spre propria lor pieire.

\VS{17}Așadar, voi, iubiților, cunoscând mai dinainte aceste lucruri, păziți-vă, ca nu cumva, fiind duși de rătăcirea celor răi, să cădeți din propria voastră statornicie.

\VS{18}Ci creșteți în harul și în cunoașterea Domnului și Mântuitorului nostru Isus Hristos. A lui să fie gloria acum și în veci. Amin.


\par }